\documentclass[11pt]{article}
\usepackage[a4paper,margin=27mm]{geometry}
\usepackage{fontspec,unicode-math}
\setmainfont{STIXGeneral}
\setmathfont{STIX Two Math}
\setmonofont{Menlo}
\usepackage{amsmath,amsthm,booktabs,hyperref}
\newtheorem{theorem}{Theorem}
\newtheorem{proposition}[theorem]{Proposition}
\newtheorem{corollary}[theorem]{Corollary}
\newcommand{\PSL}{\operatorname{PSL}_2}
\newcommand{\PGL}{\operatorname{PGL}_2}
\newcommand{\gen}[1]{\langle #1\rangle}
\newcommand{\CG}{\mathrm C}
\newcommand{\tr}{\operatorname{tr}}
\setlength{\emergencystretch}{3em}
\title{Regular maniplexes on the projective line}
\author{human\_compression project}
\date{11 October 2026}
\begin{document}
\maketitle

\section{The source contract and common theory}
This application recovers the principal results of Leemans--Toledo
\cite{LT} from the common paper \cite{U}. The source is the eleven-page
author manuscript dated 15 May 2023. The claims recovered here are its
rank restriction, Theorem~3.10, Proposition~4.2, Theorem~4.4, the two
infinite constructions and Corollaries~4.6 and~4.8. We also correct the
source's restricted list of possible face counts and several auxiliary
statements. The corrections are distinguished from the source's results;
no priority or new conjecture claim is made.

Write $L(q)=\PSL(q)$, $P(q)=\PGL(q)$, and
$G_I=\gen{r_i:i\in I}$ for a marked group. A string marking is an ordered
list of distinct nonidentity involutions generating $G$, with
$r_i r_j=r_j r_i$ for $|i-j|>1$. Its type lists the orders of adjacent
products. The intersection property is not part of this definition.
For rank four put $H=G_{\{0,1,2\}}$ and $K=G_{\{1,2,3\}}$.

We use the common paper's projective string theorem, two-proper-end
classification, geometric identification and half-turn extension theorem
\cite[Theorems~4, 6, 8, 9 and Corollary~7]{U}. These are proved there for string markings,
so their use here requires no intersection-property assumption. In
particular, the local $S_4/A_5$ classification is the exhaustive
$432$-triple computation in the common required program: generation and
commuting endpoints are admission conditions, and the intersection
property is verified afterward. The same program contains the complete
local presentations, ten universal groups and literal correction witnesses
used below. Its full cost is part of the common theory.

The outside inputs retained here are the usual projective subgroup and
torus theory declared in the common paper, Giudici's precise maximal
$A_5$ and $S_4$ cases \cite[Theorem~2.2(6),(8)]{Giudici}, and Dirichlet's
theorem for primes in coprime progressions. The final Suzuki observation
uses the stated natural-action fact. We do not import an earlier version
of this application's proof as an additional premise.

\section{From marked groups to maniplexes}
For completeness we give the category interface. A maniplex has flags and
fixed-point-free involutory adjacencies $s_i$, such that $s_i s_j$ is
fixed-point-free for $i\ne j$, nonconsecutive adjacencies commute, and the
adjacencies act transitively. An automorphism preserves every adjacency;
a maniplex is regular when its automorphism group is transitive on flags.
As in \cite[Remark~2.1]{LT}, a rank-three maniplex is called a map; this
does not assert that its face poset is an abstract polytope.

Given a string marking of $G$, take its elements as flags and let $s_i$
act by left multiplication by $r_i$. Nonidentity gives fixed-point-free
adjacencies, distinctness gives fixed-point-free products of two different
adjacencies, and generation gives connectedness. An adjacency-preserving
permutation $f$ commutes with all left multiplications. If $f(1)=h$, then
$f(g)=gh$ for every $g\in G$; conversely every such right multiplication
preserves the adjacencies. Thus the full colour-preserving automorphism
group is the right regular group, isomorphic to $G$ by taking inverse
right multiplications. This proves equality with the full automorphism
group, rather than merely the presence of a flag-transitive subgroup.

Conversely, in a connected maniplex an automorphism fixing one flag fixes
all flags, by following coloured paths. Hence in a regular maniplex there
is a unique automorphism taking a chosen base flag to each adjacent flag.
Call these automorphisms $r_i$. They commute with every adjacency. Moving
the base flag along a path shows that its terminal automorphism is the
reversed product of the corresponding $r_i$; these therefore generate.
The fixed-point-free and string conditions give distinct nonidentity
involutions and the required nonconsecutive commutations. Changing the
base flag conjugates the marking. This proves the correspondence up to
marking-preserving isomorphism.

An $i$-face is an orbit of all adjacencies except $i$. In the regular model
its number is the index of the corresponding generated subgroup. Reversal
of the marking gives the dual maniplex, so an automorphism sending every
$r_i$ to $r_{3-i}$ is a duality. The subgroup of even words has index at
most two. For nonabelian simple $G$ it must be $G$: otherwise $G$ would
have a nontrivial quotient of order at most two. The even adjacencies then
act transitively, which is exactly nonorientability in the source's
convention.

\section{Classification and extension as corollaries}
\begin{corollary}[Leemans--Toledo]\label{cor:classification}
If $L(q)$ admits a string marking of rank at least four, then $q$ is odd
and the rank is four. Each end group is $S_4$, $A_5$ or $L(q)$.
If both ends are proper, the marked object is uniquely the $11$-cell for
$q=11$, of type $\{3,5,3\}$, or the $57$-cell for $q=19$, of type
$\{5,3,5\}$. Both satisfy the intersection property and are self-dual.
\end{corollary}
\begin{proof}
Apply the projective string theorem and two-proper-end theorem of
\cite[Theorems~4 and~6]{U} to $G=L(q)$. The same-field $\PSL(q)$ alternative in the
proper-end list for $\PGL(q)$ is absent for this group. The finite
theorem leaves $L(11)$, $L(19)$ or $S_5$; simplicity excludes $S_5$.
The strictly increasing odd-field order $q(q^2-1)/2$ identifies the
remaining field sizes as $11$ and $19$. The common geometric theorem
identifies the two surviving marked groups and their
dualities. Its subgroup indices give $11$ vertices and $11$ facets in
the first case, and $57$ of each in the second. Their facets are the
hemi-icosahedron and hemi-dodecahedron, respectively, and their vertex
figures are the dual local maps. This is precisely the source's
Theorem~3.10, Proposition~4.2 and characteristic-two exclusion, together
with the asserted rank bound. The small groups $L(2)$ and $L(3)$ are
included in the common theorem's explicit exceptions check.
\end{proof}

\begin{corollary}[Leemans--Toledo, Theorem~4.4]\label{cor:extension}
Let $q$ be odd, choose $\alpha\in\{1,-1\}$ with
$q\equiv\alpha\pmod4$, and set $N=(q-\alpha)/2$. A rank-three string
marking $(r_0,r_1,r_2)$ of $L(q)$ whose first adjacent order is an odd
divisor of $N$ extends to rank four with the same group and the original
map as its unique facet.
\end{corollary}
\begin{proof}
The common indecomposability argument makes both old adjacent labels
at least three. The full-end half-turn extension theorem of
\cite[Theorem~9]{U} supplies a new
involution $r_3$ and proves all four generators distinct, their string
commutations and generation. It preserves the first three generators,
which generate $L(q)$. The preceding category interface therefore gives
one facet, with exactly the original marked map. This uses the full
odd-divisor range of the common theorem, not a finite list of fields.
\end{proof}

The same common extension theorem has a maximal-end form. If
$H<L(q)$ is maximal, its first adjacent order is an odd divisor of $N$,
and $\CG_H(r_0,r_1)=1$, the torus half-turn lies outside $H$ and the
extended marking generates $L(q)$. Its actual first-three subgroup is
$H$. We now supply the source-specific inputs for both constructions.

\section{The source-marked family with two full ends}
\begin{theorem}[Leemans--Toledo, Lemma~4.5 and Corollary~4.6]
\label{thm:full-family}
For every prime $p\equiv1\pmod{12}$ there is a rank-four regular
nonorientable maniplex with group $L(p)$ and both end groups equal to
$L(p)$. Its first three generators have type $\{3,p\}$ and give the
source's marked map. There are infinitely many such maniplexes.
\end{theorem}
\begin{proof}
Consider the projective classes of
\[
 R_0=\begin{pmatrix}0&1\\1&0\end{pmatrix},\quad
 R_1=\begin{pmatrix}1&0\\1&-1\end{pmatrix},\quad
 R_2=\begin{pmatrix}-1&0\\0&1\end{pmatrix},\quad
 R_3=\begin{pmatrix}1&-2\\2&-1\end{pmatrix}.
 \tag{1}\label{eq:matrices}
\]
The first three determinants are $-1$ and the last is $3$.
Since $4\mid p-1$, $-1$ is a square. Choose $\zeta$ of order three;
then $(2\zeta+1)^2=-3$. Multiplying by a square root of $-1$ shows that
$3$ is a square too. Thus all four projective classes lie in $L(p)$.
Their matrices have trace zero and nonzero determinant, and none is
scalar, so they are nonidentity involutions.

Direct multiplication gives $r_0r_2=r_2r_0$,
$|r_0r_1|=3$, and
\[
 [R_1R_2]=\left[\begin{pmatrix}1&0\\1&1\end{pmatrix}\right].
 \tag{2}\label{eq:unipotent}
\]
The latter has exact order $p$. The adjacent orders and the displayed
matrices also show that the first three classes are distinct. Their
powers supply all lower unipotent elementary matrices over $\mathbb F_p$;
conjugation by $R_0$ supplies all upper ones. These generate
$\operatorname{SL}_2(p)$ by elementary row elimination. To include the
last diagonal step explicitly, let
\[
 u(t)=\begin{pmatrix}1&t\\0&1\end{pmatrix},\quad
 v(t)=\begin{pmatrix}1&0\\t&1\end{pmatrix},\quad
 h(t)=u(t)v(-t^{-1})u(t)=\begin{pmatrix}0&t\\-t^{-1}&0\end{pmatrix}.
\]
For $t\ne0$ one has $h(t)h(-1)=\operatorname{diag}(t,t^{-1})$.
The matrix $h(1)$ permits a row interchange when required, and the upper
and lower matrices eliminate the other entries. Hence the first three
projective classes generate $L(p)$.

Here $3\mid(p-1)/2$, so the common extension theorem applies.
Put $A=R_0R_1$. The identity $2A-\tr(A)I=R_3$, together with the
common dihedral-centralizer formula in \cite[Lemma~1]{U}, proves that $r_3$
centralizes $r_0,r_1$. Both old adjacent pairs are noncommuting, so the
full-end extension theorem also proves that this fourth involution
repeats none of the old generators. Equivalently, it is the unique
torus half-turn in Corollary~\ref{cor:extension}. The other end contains
$r_1r_2$ of order $p\ge13$. Corollary~\ref{cor:classification} makes it
$S_4$, $A_5$ or $L(p)$, and the first two have no element of order $p$.
Both ends are therefore full.

The construction agrees with the actual source marking, not merely its
group order. Choose $a\in\mathbb F_p$ with $a^2=-1$. The three source
matrices are
\[
 S_0=\begin{pmatrix}0&1\\-1&0\end{pmatrix},\quad
 S_1=\begin{pmatrix}a&0\\1&-a\end{pmatrix},\quad
 S_2=\begin{pmatrix}-a&0\\0&a\end{pmatrix}.
\]
For $C=\operatorname{diag}(-a,1)$ their conjugates are respectively
$-aR_0,aR_1,aR_2$. Thus the ordered projective triples are conjugate.
Uniqueness of the nonidentity common-centralizer element identifies the
fourth marker as well. The group--maniplex correspondence gives one
vertex and one facet. Simplicity proves nonorientability. Finally,
Dirichlet's theorem supplies infinitely many primes in the progression,
and the strictly increasing orders $p(p^2-1)/2$ distinguish the groups.
\end{proof}

\section{Proper alternating and symmetric ends}
\begin{theorem}[Leemans--Toledo, Proposition~4.7 and Corollary~4.8]
\label{thm:A5}
For every prime $p\equiv\pm1\pmod{20}$, $p\ne19$, there is a rank-four
regular nonorientable maniplex with group $L(p)$, one vertex and
$|L(p)|/60$ hemi-dodecahedral facets. Its end groups are $L(p)$ and
$A_5$. There are infinitely many such maniplexes.
\end{theorem}
\begin{proof}
Giudici's retained theorem gives a maximal $A_5$ subgroup $H$ for every
prime in the stated range. Choose the common local marking of type $\{5,3\}$
with Petrie order five. Its first pair generates $D_{10}$.
The centralizer of a $5$-cycle in $A_5$ is its cyclic group of order five;
the dihedral reflection inverts it, fixing only its identity. Thus
$\CG_H(r_0,r_1)=1$.

Let $p\equiv\alpha\pmod4$, $\alpha\in\{1,-1\}$. If
$p\equiv1\pmod{20}$ then $\alpha=1$ and $10\mid(p-1)/2$; if
$p\equiv-1\pmod{20}$ then $\alpha=-1$ and $10\mid(p+1)/2$.
The first adjacent order five is therefore an odd divisor of the even
torus order $N=(p-\alpha)/2$. The maximal-end extension theorem gives
a fourth involution outside $H$, and the resulting marking generates
$L(p)$ with first-three subgroup $H$.

If its opposite end were proper, Corollary~\ref{cor:classification}
would give $p=11$ or $19$. Neither is among the present parameters.
The opposite end is full, giving the asserted face counts. The marking
of $H$ identifies the facets as hemi-dodecahedral. At $p=19$ the same
extension construction still makes a string marking, but this argument
does not force the opposite end to be full; this is why the source
excludes that value from its family assertion. Simplicity and Dirichlet
complete the nonorientability and infinitude claims.
\end{proof}

\begin{theorem}[Correction: the omitted symmetric-end family]
\label{thm:S4}
For every prime $p\equiv\pm1\pmod{24}$ there is a rank-four regular
nonorientable maniplex with group $L(p)$, one vertex and $|L(p)|/24$
facets. Its facet marking has type $\{3,4\}$ and group $S_4$, and its
other end group is $L(p)$. There are infinitely many such examples.
\end{theorem}
\begin{proof}
The retained maximal-subgroup theorem gives a maximal $S_4$ subgroup
$H$ for these primes. Use its common local marking of type $\{3,4\}$
with Petrie order three. The first pair generates an $S_3$ fixing one
point in the natural degree-four action. An element centralizing this
$S_3$ must preserve that fixed point; on the other three points it lies
in the centre of their full symmetric group and is therefore trivial.
Hence $\CG_H(r_0,r_1)=1$.

For $p\equiv1\pmod{24}$, $N=(p-1)/2$ is divisible by twelve; for
$p\equiv-1\pmod{24}$, $N=(p+1)/2$ is divisible by twelve.
The first adjacent order three therefore meets the maximal-end
extension hypothesis. The resulting fourth involution lies outside $H$
and generates $L(p)$ with it. These primes satisfy $p\ge23$, so the
opposite end cannot be proper by Corollary~\ref{cor:classification}.
Its fullness gives the face counts; simplicity gives nonorientability,
and Dirichlet gives infinitude.
\end{proof}

Consequently the end-group and face-count alternatives, up to duality,
are exactly
\[
\begin{array}{c|c}
\text{end groups}&(\text{vertices},\text{facets})\\\hline
G,G&(1,1)\\
G,S_4&(1,|G|/24)\\
G,A_5&(1,|G|/60)\\
\text{two proper ends}&(11,11)\text{ or }(57,57).
\end{array}
\]
This lists the possibilities for end groups and counts, not all marked
isomorphism classes. Both one-proper-end cases are non-self-dual because
the two counts differ. The full/full case can be self-dual.

\section{Literal corrections and exact witnesses}
The source's introduction and conclusion omit the $S_4$ alternative,
although its Theorem~3.10 and Remark~4.1 permit it.
Theorem~\ref{thm:S4} gives a uniform correction. Independently of that
uniform argument, the common program supplies a literal degree-$24$
marking of $L(23)$ of type $\{3,4,12\}$, with whole order $6072$ and
end orders $24,6072$. It checks distinct involutions, all distant
commutations, generation and the $S_4$ identification. This gives one
vertex and $253$ facets directly.

For the source page~9 self-duality sentence, use the following marking
of $L(13)$ in degree fourteen:
\[
\begin{aligned}
r_0&=(3\ 9)(4\ 10)(5\ 11)(6\ 12)(7\ 13)(8\ 14),\\
r_1&=(1\ 4)(2\ 14)(3\ 9)(5\ 13)(6\ 11)(7\ 12),\\
r_2&=(1\ 2)(3\ 7)(4\ 6)(8\ 14)(9\ 13)(10\ 12),\\
r_3&=(1\ 2)(4\ 14)(5\ 13)(6\ 12)(7\ 11)(8\ 10).
\end{aligned}
\]
The common program verifies whole and end orders $1092$, type
$\{3,7,3\}$, and all four conjugacies $r_i^d=r_{3-i}$ for
\[
 d=(1\ 3)(2\ 9)(4\ 7)(5\ 8)(6\ 12)(10\ 11)(13\ 14).
\]
Thus conjugation by $d$ induces a reversal automorphism, and the
corresponding full/full maniplex is self-dual. The elementary unequal
face-count obstruction applies to the one-proper-end cases instead.

The source's Lemma~3.3 asserts involution-centralizer maximality too
broadly. The common program verifies in $L(5)\cong A_5$ that such a
centralizer is a four-group properly contained in $A_4$. Our argument
uses the common-centralizer theorem and torus normalizers, without
assuming maximality of every involution centralizer. The broad
uniqueness assertions in source Lemma~3.4 are not retained inputs.
Likewise, Lemmas~3.7 and~3.8 must concern natural subfield embeddings,
not arbitrary abstract isomorphisms: $A_5\cong L(5)$ and
$S_4\cong P(3)$ occur as allowed proper ends. Finally, the source first
correctly calls the local map of type $\{5,3\}$ hemi-dodecahedral,
then calls its facets hemi-icosahedral. The actual marking has been
preserved here and the name made consistent.

\section{Remaining scope and verification}
The source's Suzuki observation follows from its retained natural-action
fact: every involution in a simple Suzuki group fixes a unique point,
and commuting involutions have the same fixed point. At rank four,
$r_0,r_2,r_3$ must fix the same point, and then so does $r_1$ because
it commutes with $r_3$. Every later generator commutes with $r_0$ and
fixes that point as well. This contradicts generation of the transitive
group. The nonsimple group $\operatorname{Sz}(2)$ cannot be generated
by involutions, since its quotient $C_4$ cannot be so generated.
This observation is included with its original attribution; no additional
Suzuki classification is claimed.

The final PGL remarks and four-family conjecture in \cite[Section~5]{LT}
are a separate contract. No complete enumeration of rank-four PSL
markings is claimed here. Unused auxiliary subgroup multiplicities and
uniqueness clauses are not presented as separately recovered theorems.
These qualifications do not change the complete principal contract
stated in Section~1.

The common finite program is a required mathematical component. Its
exact finite-group and finitely presented-group algorithms remain an
explicit GAP trust boundary. Finite samples of the positive families
are useful diagnostic checks; the proofs above establish their full
prime ranges and do not infer them from samples. The source basis
adapter is proved by the displayed conjugacy for every admissible prime.
The generic extension now has a separate Lean certificate proving
actual tuple injectivity, string relations, generation and preserved
prefix under its stated hypotheses. Earlier Lean certificates prove
scoped centralizer and rank interfaces. They do not formalize torus
existence, Dickson or Giudici classification, the finite universal table,
Dirichlet, the maniplex correspondence or every theorem of this paper.

\begin{thebibliography}{9}
\bibitem{LT} D.~Leemans and M.~Toledo,
\emph{Maniplexes with automorphism group $\operatorname{PSL}(2,q)$},
author manuscript dated 15 May 2023, eleven pages.
\bibitem{U} human\_compression project,
\emph{Common theory for projective markings and rotations},
common Topic E theory, accompanying \texttt{umbrella.tex} and complete
\texttt{required\_finite.g}. Precise theorem labels and final edition
hashes are bound in the joint proof map.
\bibitem{Giudici} M.~Giudici,
\emph{Maximal subgroups of almost simple groups with socle
$\operatorname{PSL}(2,q)$}, arXiv:math/0703685v1, 2007,
Theorem~2.2, physical and printed page~4.
\end{thebibliography}
\end{document}
